% Two ways to build a 4x4x4 matrix multiply-accumulate unit from the
% dot-product step of figure l12-dot-unit.  Option 1 (left): 16 dot-product
% units of 4 FMAs each, one per output element, all working in the same
% cycle (64 FMAs).  Option 2 (right): fewer units, pipelined; the row-column
% pairs (r1,c1), (r1,c2), (r1,c3), ... pass through load, compute and store
% stages one cycle apart.
% Redraws the outline figure m12l3_matrix.
% Usage: % Two ways to build a 4x4x4 matrix multiply-accumulate unit from the
% dot-product step of figure l12-dot-unit.  Option 1 (left): 16 dot-product
% units of 4 FMAs each, one per output element, all working in the same
% cycle (64 FMAs).  Option 2 (right): fewer units, pipelined; the row-column
% pairs (r1,c1), (r1,c2), (r1,c3), ... pass through load, compute and store
% stages one cycle apart.
% Redraws the outline figure m12l3_matrix.
% Usage: % Two ways to build a 4x4x4 matrix multiply-accumulate unit from the
% dot-product step of figure l12-dot-unit.  Option 1 (left): 16 dot-product
% units of 4 FMAs each, one per output element, all working in the same
% cycle (64 FMAs).  Option 2 (right): fewer units, pipelined; the row-column
% pairs (r1,c1), (r1,c2), (r1,c3), ... pass through load, compute and store
% stages one cycle apart.
% Redraws the outline figure m12l3_matrix.
% Usage: % Two ways to build a 4x4x4 matrix multiply-accumulate unit from the
% dot-product step of figure l12-dot-unit.  Option 1 (left): 16 dot-product
% units of 4 FMAs each, one per output element, all working in the same
% cycle (64 FMAs).  Option 2 (right): fewer units, pipelined; the row-column
% pairs (r1,c1), (r1,c2), (r1,c3), ... pass through load, compute and store
% stages one cycle apart.
% Redraws the outline figure m12l3_matrix.
% Usage: \input{figures/l12-mma-options}
% Width: body (120 mm); the drawing is about 114 mm wide.
\begin{tikzpicture}[font=\scriptsize\sffamily,
    du/.style={draw, fill=ln-accent!12, minimum width=8mm,
      minimum height=5.5mm, inner sep=0pt},
    st/.style={draw, minimum width=9mm, minimum height=4.6mm, inner sep=0pt,
      text height=1.6ex, text depth=0.4ex}]
  % option 1
  \node[anchor=west] at (0,8mm)
    {\iflangja{方式1：並列ユニット}{Option 1: parallel units}};
  \foreach \i in {1,...,4} {
    \foreach \j in {1,...,4} {
      \node[du] at ({(\j-1)*9.5mm+4mm},{-(\i-1)*7mm}) {$c_{\i\j}$};
    }
  }
  \node[anchor=north west, align=left, text=ln-muted] at (0,-24mm)
    {\iflangja{4 FMA × 16 要素 = 64 FMA/サイクル}{4 FMAs $\times$ 16 elements = 64 FMAs per cycle}};
  % option 2
  \begin{scope}[xshift=46mm]
    \node[anchor=west] at (0,8mm)
      {\iflangja{方式2：パイプライン化}{Option 2: pipelined unit}};
    \foreach \r/\lab in {0/{(r1,c1)},1/{(r1,c2)},2/{(r1,c3)},3/{(r1,c4)}} {
      \node[anchor=east] at (13mm,{-\r*6mm}) {\lab};
      \node[st, fill=ln-tip!20] at ({17.5mm+\r*9mm},{-\r*6mm})
        {\iflangja{ロード}{load}};
      \node[st, fill=ln-accent!18] at ({26.5mm+\r*9mm},{-\r*6mm})
        {\iflangja{演算}{comp}};
      \node[st, fill=ln-caution!22] at ({35.5mm+\r*9mm},{-\r*6mm})
        {\iflangja{ストア}{store}};
    }
    \draw[->, ln-muted] (13mm,-24mm) -- node[below] {\iflangja{サイクル}{cycle}} (68mm,-24mm);
  \end{scope}
\end{tikzpicture}

% Width: body (120 mm); the drawing is about 114 mm wide.
\begin{tikzpicture}[font=\scriptsize\sffamily,
    du/.style={draw, fill=ln-accent!12, minimum width=8mm,
      minimum height=5.5mm, inner sep=0pt},
    st/.style={draw, minimum width=9mm, minimum height=4.6mm, inner sep=0pt,
      text height=1.6ex, text depth=0.4ex}]
  % option 1
  \node[anchor=west] at (0,8mm)
    {\iflangja{方式1：並列ユニット}{Option 1: parallel units}};
  \foreach \i in {1,...,4} {
    \foreach \j in {1,...,4} {
      \node[du] at ({(\j-1)*9.5mm+4mm},{-(\i-1)*7mm}) {$c_{\i\j}$};
    }
  }
  \node[anchor=north west, align=left, text=ln-muted] at (0,-24mm)
    {\iflangja{4 FMA × 16 要素 = 64 FMA/サイクル}{4 FMAs $\times$ 16 elements = 64 FMAs per cycle}};
  % option 2
  \begin{scope}[xshift=46mm]
    \node[anchor=west] at (0,8mm)
      {\iflangja{方式2：パイプライン化}{Option 2: pipelined unit}};
    \foreach \r/\lab in {0/{(r1,c1)},1/{(r1,c2)},2/{(r1,c3)},3/{(r1,c4)}} {
      \node[anchor=east] at (13mm,{-\r*6mm}) {\lab};
      \node[st, fill=ln-tip!20] at ({17.5mm+\r*9mm},{-\r*6mm})
        {\iflangja{ロード}{load}};
      \node[st, fill=ln-accent!18] at ({26.5mm+\r*9mm},{-\r*6mm})
        {\iflangja{演算}{comp}};
      \node[st, fill=ln-caution!22] at ({35.5mm+\r*9mm},{-\r*6mm})
        {\iflangja{ストア}{store}};
    }
    \draw[->, ln-muted] (13mm,-24mm) -- node[below] {\iflangja{サイクル}{cycle}} (68mm,-24mm);
  \end{scope}
\end{tikzpicture}

% Width: body (120 mm); the drawing is about 114 mm wide.
\begin{tikzpicture}[font=\scriptsize\sffamily,
    du/.style={draw, fill=ln-accent!12, minimum width=8mm,
      minimum height=5.5mm, inner sep=0pt},
    st/.style={draw, minimum width=9mm, minimum height=4.6mm, inner sep=0pt,
      text height=1.6ex, text depth=0.4ex}]
  % option 1
  \node[anchor=west] at (0,8mm)
    {\iflangja{方式1：並列ユニット}{Option 1: parallel units}};
  \foreach \i in {1,...,4} {
    \foreach \j in {1,...,4} {
      \node[du] at ({(\j-1)*9.5mm+4mm},{-(\i-1)*7mm}) {$c_{\i\j}$};
    }
  }
  \node[anchor=north west, align=left, text=ln-muted] at (0,-24mm)
    {\iflangja{4 FMA × 16 要素 = 64 FMA/サイクル}{4 FMAs $\times$ 16 elements = 64 FMAs per cycle}};
  % option 2
  \begin{scope}[xshift=46mm]
    \node[anchor=west] at (0,8mm)
      {\iflangja{方式2：パイプライン化}{Option 2: pipelined unit}};
    \foreach \r/\lab in {0/{(r1,c1)},1/{(r1,c2)},2/{(r1,c3)},3/{(r1,c4)}} {
      \node[anchor=east] at (13mm,{-\r*6mm}) {\lab};
      \node[st, fill=ln-tip!20] at ({17.5mm+\r*9mm},{-\r*6mm})
        {\iflangja{ロード}{load}};
      \node[st, fill=ln-accent!18] at ({26.5mm+\r*9mm},{-\r*6mm})
        {\iflangja{演算}{comp}};
      \node[st, fill=ln-caution!22] at ({35.5mm+\r*9mm},{-\r*6mm})
        {\iflangja{ストア}{store}};
    }
    \draw[->, ln-muted] (13mm,-24mm) -- node[below] {\iflangja{サイクル}{cycle}} (68mm,-24mm);
  \end{scope}
\end{tikzpicture}

% Width: body (120 mm); the drawing is about 114 mm wide.
\begin{tikzpicture}[font=\scriptsize\sffamily,
    du/.style={draw, fill=ln-accent!12, minimum width=8mm,
      minimum height=5.5mm, inner sep=0pt},
    st/.style={draw, minimum width=9mm, minimum height=4.6mm, inner sep=0pt,
      text height=1.6ex, text depth=0.4ex}]
  % option 1
  \node[anchor=west] at (0,8mm)
    {\iflangja{方式1：並列ユニット}{Option 1: parallel units}};
  \foreach \i in {1,...,4} {
    \foreach \j in {1,...,4} {
      \node[du] at ({(\j-1)*9.5mm+4mm},{-(\i-1)*7mm}) {$c_{\i\j}$};
    }
  }
  \node[anchor=north west, align=left, text=ln-muted] at (0,-24mm)
    {\iflangja{4 FMA × 16 要素 = 64 FMA/サイクル}{4 FMAs $\times$ 16 elements = 64 FMAs per cycle}};
  % option 2
  \begin{scope}[xshift=46mm]
    \node[anchor=west] at (0,8mm)
      {\iflangja{方式2：パイプライン化}{Option 2: pipelined unit}};
    \foreach \r/\lab in {0/{(r1,c1)},1/{(r1,c2)},2/{(r1,c3)},3/{(r1,c4)}} {
      \node[anchor=east] at (13mm,{-\r*6mm}) {\lab};
      \node[st, fill=ln-tip!20] at ({17.5mm+\r*9mm},{-\r*6mm})
        {\iflangja{ロード}{load}};
      \node[st, fill=ln-accent!18] at ({26.5mm+\r*9mm},{-\r*6mm})
        {\iflangja{演算}{comp}};
      \node[st, fill=ln-caution!22] at ({35.5mm+\r*9mm},{-\r*6mm})
        {\iflangja{ストア}{store}};
    }
    \draw[->, ln-muted] (13mm,-24mm) -- node[below] {\iflangja{サイクル}{cycle}} (68mm,-24mm);
  \end{scope}
\end{tikzpicture}
